%%%%%%%%%%%%%%%%%%%%%%%%%%%%%%%%%%%%%%%%%%%%%%%%%%%%%%%%%%%%%%%%%%%%%%%%%%%%%%%%%%
\begin{frame}[fragile]\frametitle{}
\begin{center}
{\Large The 5G Framework}

\vspace{0.3cm}
{\em Grids, Groups, Graphs, Geodesics, Gauges}
\end{center}
\end{frame}

%%%%%%%%%%%%%%%%%%%%%%%%%%%%%%%%%%%%%%%%%%%%%%%%%%%%%%%%%%%
\begin{frame}[fragile]\frametitle{Why Do We Need So Many Architectures?}

\begin{center}
\textbf{Different data structures demand different neural network designs.}
\end{center}

\begin{center}
\begin{tabular}{llll}
\hline
\textbf{Data Type} & \textbf{Structure} & \textbf{Key Property} & \textbf{Architecture} \\
\hline
Tabular & Fixed vectors & Independence & MLP \\
Images & Regular grid & Translation & CNN \\
Sequences & 1D chain & Order & RNN \\
Sets & Unordered & Permutation & Attention \\
Graphs & Irregular & Connectivity & GNN \\
\hline
\end{tabular}
\end{center}

\vspace{0.3cm}

\begin{block}{The Geometric Deep Learning Hypothesis}
\begin{center}
\textbf{The success of a neural architecture is determined by how well its geometric biases match the geometric structure of the data.}
\end{center}
\end{block}

\end{frame}

%%%%%%%%%%%%%%%%%%%%%%%%%%%%%%%%%%%%%%%%%%%%%%%%%%%%%%%%%%%%%%%%%%%%%%%%%%%%%%%%%%
\begin{frame}[fragile]\frametitle{Felix Klein's Vision Meets Deep Learning}

\begin{columns}
\begin{column}{0.6\textwidth}
\textbf{Klein's Erlangen Programme (1872):}
\begin{itemize}
\item Geometry = study of properties that are \textbf{invariant under a symmetry group}
\item Unified all geometries through group theory
\end{itemize}

\vspace{0.4cm}
\textbf{Modern Translation for Neural Networks:}
\begin{itemize}
\item Architecture = set of computations \textbf{invariant under a group of transformations}
\item Unified framework for all data domains
\item Same mathematical idea, 150 years later
\end{itemize}
\end{column}
\begin{column}{0.4\textwidth}
\begin{center}
\textit{[image: Felix Klein / Erlangen Programme diagram]}

\vspace{0.4cm}
{\small ``Geometry is the study of invariant properties under group actions''}
\end{center}
\end{column}
\end{columns}

\end{frame}

%%%%%%%%%%%%%%%%%%%%%%%%%%%%%%%%%%%%%%%%%%%%%%%%%%%%%%%%%%%%%%%%%%%%%%%%%%%%%%%%%%
\begin{frame}[fragile]\frametitle{The 5G Framework: A Tourist Map}

\begin{block}{Five geometric domains, increasing in generality}
\begin{itemize}
  \item \textbf{Grid}: Regular lattices (images, video) $\rightarrow$ \textbf{Translation equivariance} $\rightarrow$ CNNs
  \item \textbf{Group}: Homogeneous spaces (sphere, rotation group) $\rightarrow$ \textbf{Group equivariance} $\rightarrow$ Spherical CNNs
  \item \textbf{Graph}: Discrete irregular structures $\rightarrow$ \textbf{Permutation invariance} $\rightarrow$ GNNs
  \item \textbf{Geodesic}: Continuous manifolds (3D meshes, brain cortex) $\rightarrow$ \textbf{Diffeomorphism invariance} $\rightarrow$ Mesh CNNs
  \item \textbf{Gauge}: Vector bundles (weather fields, EM fields) $\rightarrow$ \textbf{Gauge equivariance} $\rightarrow$ Gauge CNNs
\end{itemize}
\end{block}

\vspace{0.2cm}
Each ``G'' is more general than the previous --- Grid is a special case of Group, which is a special case of Graph, and so on.

\end{frame}

%%%%%%%%%%%%%%%%%%%%%%%%%%%%%%%%%%%%%%%%%%%%%%%%%%%%%%%%%%%%%%%%%%%%%%%%%%%%%%%%%%
\begin{frame}[fragile]\frametitle{Invariance vs.\ Equivariance: What Is the Difference?}

\begin{columns}
\begin{column}{0.5\textwidth}
\textbf{Invariance:} Output does \emph{not} change.
$$f(g \cdot x) = f(x)$$

\textbf{Example:} ``Is this a cat?'' answer is the same wherever the cat is in the image.

\vspace{0.4cm}
\textbf{Equivariance:} Output transforms \emph{predictably}.
$$f(g \cdot x) = \rho(g) \cdot f(x)$$

\textbf{Example:} If the image is translated, the feature map is also translated by the same amount.
\end{column}
\begin{column}{0.5\textwidth}
\begin{exampleblock}{CNN Translation Equivariance}
\texttt{conv(translate(image))}\\
\texttt{= translate(conv(image))}
\end{exampleblock}

\vspace{0.3cm}
\begin{alertblock}{Why Equivariance $>$ Invariance?}
Equivariance preserves \textbf{spatial relationships}, enabling hierarchical processing. Global invariance is only needed at the very last layer.
\end{alertblock}
\end{column}
\end{columns}

\end{frame}

%%%%%%%%%%%%%%%%%%%%%%%%%%%%%%%%%%%%%%%%%%%%%%%%%%%%%%%%%%%%%%%%%%%%%%%%%%%%%%%%%%
\begin{frame}[fragile]\frametitle{The Permutation Challenge in Graphs}

\begin{columns}
\begin{column}{0.4\textwidth}
\textbf{All orderings of $\{A,B,C\}$:}
\begin{itemize}
  \item $\{A, B, C\}$
  \item $\{A, C, B\}$
  \item $\{B, A, C\}$
  \item $\{B, C, A\}$
  \item $\{C, A, B\}$
  \item $\{C, B, A\}$
\end{itemize}
Total: $3! = 6$
\end{column}
\begin{column}{0.6\textwidth}
\textbf{Naive approach:} learn for each ordering separately $\rightarrow$ 6$\times$ more data, no generalisation.

\vspace{0.3cm}
\textbf{Geometric ML approach:} learn a \textbf{permutation-invariant} function --- works for all orderings with the same parameters.

\begin{block}{Mathematical Formulation}
$$f(\pi \cdot X) = f(X) \quad \forall \pi \in S_n$$
where $\pi$ is any permutation, $S_n$ is the symmetric group.
\end{block}
\end{column}
\end{columns}

\end{frame}

%%%%%%%%%%%%%%%%%%%%%%%%%%%%%%%%%%%%%%%%%%%%%%%%%%%%%%%%%%%%%%%%%%%%%%%%%%%%%%%%%%
\begin{frame}[fragile]\frametitle{Grid Domain: CNNs and Translation Equivariance}

\begin{columns}
\begin{column}{0.6\textwidth}
\textbf{Key properties of the grid domain:}
\begin{itemize}
\item Fixed, regular neighbourhood structure
\item Translation symmetry
\item Local connectivity
\end{itemize}

\vspace{0.3cm}
\textbf{CNN Convolution:}
$$(\mathbf{f} * \mathbf{g})[i,j] = \sum_{m,n} \mathbf{f}[m,n] \cdot \mathbf{g}[i-m, j-n]$$

\textbf{Equivariance:}
$$\text{Conv}(T_v(\mathbf{x})) = T_v(\text{Conv}(\mathbf{x}))$$
\end{column}
\begin{column}{0.4\textwidth}
\begin{center}
\textit{[image: CNN sliding window on a regular image grid]}

\vspace{0.3cm}
{\small Convolution preserves spatial relationships}
\end{center}
\end{column}
\end{columns}

\begin{alertblock}{Limitation}
CNNs fail on irregular, non-grid data (graphs, point clouds, 3D meshes).
\end{alertblock}

\end{frame}

%%%%%%%%%%%%%%%%%%%%%%%%%%%%%%%%%%%%%%%%%%%%%%%%%%%%%%%%%%%%%%%%%%%%%%%%%%%%%%%%%%
\begin{frame}[fragile]\frametitle{Graph Domain: Permutation Invariance}

\begin{columns}
\begin{column}{0.55\textwidth}
\textbf{Key challenges:}
\begin{itemize}
\item Variable node degrees
\item No canonical node ordering
\item Irregular neighbourhood sizes
\item Arbitrary connectivity
\end{itemize}

\textbf{Graph Convolution (GCN):}
$$\mathbf{H}^{(l+1)} = \sigma\!\left(\mathbf{D}^{-\frac{1}{2}}\mathbf{A}\mathbf{D}^{-\frac{1}{2}}\mathbf{H}^{(l)}\mathbf{W}^{(l)}\right)$$

\textbf{Permutation Invariance:}
$$f(\mathbf{P}\mathbf{X},\, \mathbf{P}\mathbf{A}\mathbf{P}^T) = f(\mathbf{X},\, \mathbf{A})$$
\end{column}
\begin{column}{0.45\textwidth}
\begin{center}
\textit{[image: example graph with variable-degree nodes and irregular edges]}
\end{center}

\begin{exampleblock}{\small Message Passing}
\small
$\mathbf{m}_{ij} = M(\mathbf{h}_i, \mathbf{h}_j, \mathbf{e}_{ij})$\\
$\mathbf{h}_i' = U\!\left(\mathbf{h}_i,\, \mathrm{AGG}\!\left(\{\mathbf{m}_{ij}\}\right)\right)$
\end{exampleblock}
\end{column}
\end{columns}

\end{frame}

%%%%%%%%%%%%%%%%%%%%%%%%%%%%%%%%%%%%%%%%%%%%%%%%%%%%%%%%%%%%%%%%%%%%%%%%%%%%%%%%%%
\begin{frame}[fragile]\frametitle{The Unified Message Passing Framework}

\begin{center}
\textbf{All five geometric domains share the same message-passing structure.}
\end{center}

\begin{center}
\begin{tabular}{lll}
\hline
\textbf{Domain} & \textbf{Neighbourhood $\mathcal{N}(i)$} & \textbf{Message function $M$} \\
\hline
Grid & Fixed spatial neighbours & Convolution kernel \\
Group & Group orbit & Group convolution \\
Graph & Adjacent nodes & Learned function \\
Geodesic & Geodesic neighbours & Distance-based \\
Gauge & Parallel transport & Gauge connection \\
\hline
\end{tabular}
\end{center}

\begin{alertblock}{Universal Insight}
The 5G domains differ only in their \textbf{definition of neighbourhood} and \textbf{message function} --- the algorithmic pattern is identical.
\end{alertblock}

\end{frame}

%%%%%%%%%%%%%%%%%%%%%%%%%%%%%%%%%%%%%%%%%%%%%%%%%%%%%%%%%%%%%%%%%%%%%%%%%%%%%%%%%%
\begin{frame}[fragile]\frametitle{Spectral vs.\ Spatial Graph Convolutions}

\begin{columns}
\begin{column}{0.5\textwidth}
\textbf{Spectral approach:}
\begin{itemize}
\item Based on graph Fourier transform
\item Uses Laplacian eigenvectors
\item Global receptive field
\item Computationally expensive
\end{itemize}
$\hat{f}(\lambda_i) = \sum_{j} f(j)\, u_i(j)$
\end{column}
\begin{column}{0.5\textwidth}
\textbf{Spatial approach:}
\begin{itemize}
\item Direct neighbourhood aggregation
\item Local message passing
\item Works on unseen graphs (inductive)
\item Computationally efficient
\end{itemize}
$h_i' = \sigma\!\left(\sum_{j \in \mathcal{N}(i)} \frac{h_j W}{\sqrt{d_i d_j}}\right)$
\end{column}
\end{columns}

\vspace{0.3cm}

\begin{center}
\begin{tabular}{lccc}
\hline
\textbf{Property} & \textbf{Spectral} & \textbf{Spatial} \\
\hline
Inductive (new graphs) & No & Yes \\
Efficiency & Low & High \\
Interpretability & Mathematical & Intuitive \\
\hline
\end{tabular}
\end{center}

\end{frame}
